% Generated draft: requires mathematical and source review.
\documentclass[11pt]{article}
\usepackage[margin=25mm]{geometry}
\usepackage{fontspec}
\setmainfont{Latin Modern Roman}
\usepackage{amsmath,amssymb,mathtools,unicode-math}
\setmathfont{Latin Modern Math}
\usepackage{enumitem,needspace,xurl,hyperref}
\setlength{\parindent}{0pt}
\setlength{\parskip}{5pt}
\setlength{\emergencystretch}{3em}
\newcommand{\N}{\mathbb{N}}
\newcommand{\Z}{\mathbb{Z}}
\newcommand{\Q}{\mathbb{Q}}
\newcommand{\R}{\mathbb{R}}
\newcommand{\C}{\mathbb{C}}
\newcommand{\F}{\mathbb{F}}
\newcommand{\A}{\mathbb{A}}
\newcommand{\PP}{\mathbb{P}}
\newcommand{\E}{\mathbb{E}}
\newcommand{\eps}{\varepsilon}
\newcommand{\dd}{\,\mathrm{d}}
\begin{document}
\section*{UnsolvedMath ID: 700826}
\textbf{Source problem code:} catalogue entry
\subsection{Definitions and mathematical statement}
Let \(p\) be a prime. Here a \(p\)-group means a group in which every element has finite order equal to a power of \(p\), and a \(p\)-subgroup means a subgroup with this property. The source does not explicitly specify its convention for infinite \(p\)-groups; no finiteness, local finiteness, or bounded-exponent assumption is added here.

Let \(G\) be a finitely generated group and let \(H\leq G\) be a \(p\)-subgroup. Assume that
\[
\forall N\trianglelefteq G,\qquad N\leq H\ \Longrightarrow\ N=\{1\},
\]
and that
\[
HX=XH\qquad\text{for every subgroup }X\leq G,
\]
where \(HX=\{hx:h\in H,\ x\in X\}\) and \(XH=\{xh:x\in X,\ h\in H\}\).
Define the normal closure of \(H\) in \(G\) by
\[
H^G=\left\langle g^{-1}hg:h\in H,\ g\in G\right\rangle,
\]
and its centralizer by
\[
C_G(H^G)=\{g\in G:gy=yg\text{ for every }y\in H^G\}.
\]
Since \(H^G\) is normal in \(G\), its centralizer is normal in \(G\), so the quotient below is defined.

\textbf{Question.} For every prime \(p\) and every pair \((G,H)\) satisfying these assumptions, must
\[
G/C_G(H^G)
\]
be a \(p\)-group?
\subsection{Short English statement}
Suppose a finitely generated group \(G\) has a \(p\)-subgroup \(H\) that contains no nontrivial normal subgroup of \(G\) and satisfies \(HX=XH\) for every subgroup \(X\leq G\). Must the quotient of \(G\) by the centralizer of the normal closure of \(H\) be a \(p\)-group? The convention for infinite \(p\)-groups requires source review.
\subsection{Sources}
\begin{itemize}
\item Local UnsolvedMath registry, permanent ID \(700826\), exact source code \texttt{catalogue entry}, title \emph{Kourovka Notebook Problem 12.6}. The supplied record labels the problem \texttt{open}; this label is not an independently verified assessment.
\item G. Busetto, Problem 12.6, \emph{The Kourovka Notebook (Unsolved Problems in Group Theory)}, 12th issue (1992), as attributed in the supplied record.
\item The supplied background identifies the 21st issue edition (2026), \texttt{arXiv:1401.0300v48}. The supplied collection description names E. I. Khukhro and V. D. Mazurov as editors. No independent source consultation is asserted.
\end{itemize}
\end{document}
